\section{Analysis: Moment and Counter Policies}
\label{sec:analysis}

We write out the Adam update in a form that separates moment values from counters. Everything in this
section follows directly from the update rule of \citet{kingma2015adam}. We state it to fix precisely
what the interventions and controls do. It is not a new result.

\paragraph{Update with per-moment counters.}
Let $a_s=1-\beta_1^{s}$ and $b_s=1-\beta_2^{s}$. After the moment update with gradient $g$, let $m,v$
be the (post-gradient) moments. Standard Adam uses one counter $s$ for both moments. We allow separate
counters $s_m,s_v$, which coincide in ordinary training. The update of coordinate $i$ is
\begin{equation}
u_i \;=\; -\alpha\,\frac{m_i/a_{s_m}}{\sqrt{v_i/b_{s_v}}+\eps}
      \;=\; -\alpha\,\frac{\sqrt{b_{s_v}}}{a_{s_m}}\cdot\frac{m_i}{\sqrt{v_i}+\eps\sqrt{b_{s_v}}}.
\label{eq:update}
\end{equation}
This assumes $\eps$ is added to $\sqrt{\hat v}$, as in Algorithm~1 of \citet{kingma2015adam} and in
\texttt{torch.optim.Adam}, and not the rescaled $\hat\eps$ variant. It also assumes no weight decay (L2 or
decoupled), no AMSGrad, fixed $\beta_1,\beta_2,\alpha$, and exact arithmetic. Our implementation matches
\texttt{torch.optim.Adam} under these options (Appendix~\ref{app:checks}). Under a \emph{decoupled}
policy, a reset moment restarts its own counter at 0, so $1-\beta^{s}$ equals the EMA weight that moment
actually carries. Under \emph{shared-keep}, the counter is not changed. This is what zeroing
\texttt{exp\_avg\_sq} alone does in PyTorch.

\paragraph{Counter changes with equal moments.}
Consider two optimizer states with the same post-gradient moments $m,v$ and counters $t$ and $r$ (both
moments sharing each counter). For a coordinate with $v_i>0$, Eq.~\eqref{eq:update} gives
\begin{equation}
\frac{u^{(r)}_i}{u^{(t)}_i}
 \;=\; \kappa(r,t)\cdot\frac{\sqrt{v_i}+\eps\sqrt{b_t}}{\sqrt{v_i}+\eps\sqrt{b_r}},
\qquad
\kappa(r,t)=\frac{\sqrt{b_r}\,a_t}{a_r\sqrt{b_t}} .
\label{eq:ratio}
\end{equation}
If $\eps=0$, the ratio equals $\kappa(r,t)$ for every coordinate with $v_i>0$. Coordinates with
$v_i=0$ have received only zero gradients, so $m_i=0$ and both updates vanish. Hence, \emph{when
$\eps=0$}, changing the counter while keeping the moments multiplies the whole update by one scalar.
Our counters are shared by all tensors, so the scalar is also shared by all tensors. For $r=1$ and
$t\to\infty$, $\kappa\to\sqrt{1-\beta_2}/(1-\beta_1)$, which is $\approx0.316$ for $(0.9,0.999)$. For
finite $t$ the extra factor $a_t/\sqrt{b_t}$ applies, so $0.316$ is an approximation, not an identity.
If $\eps>0$, the second factor of Eq.~\eqref{eq:ratio} depends on $v_i$. It moves from $\kappa(r,t)$
when $\sqrt{v_i}\gg\eps$ to $a_t/a_r$ when $\sqrt{v_i}\ll\eps\sqrt{b_r}$. The scalar equivalence
therefore breaks for coordinates whose gradient scale is comparable to $\eps$
(Appendix~\ref{app:counterexample}).

\paragraph{Consequences for three arms.}
\emph{(i) \texttt{reset\_t}.} The counter reset leaves the moment recursions unchanged. If $\eps=0$,
induction over steps (equal weights $\Rightarrow$ equal gradients $\Rightarrow$ equal moments) shows that
\texttt{reset\_t} is \keep{} run with the deterministic learning-rate multiplier
$\kappa(k,t{+}k)$ at the $k$-th step after the reset. A per-tensor update-norm control that uses \keep's
own direction then reproduces \texttt{reset\_t} exactly in exact arithmetic. For this arm the control is
the same algorithm, not an independent test.
\emph{(ii) Shared-keep artifacts.} After \texttt{reset\_both}, the shared-keep and decoupled branches
start from the same zero moments and differ only in their counters ($t{+}k$ versus $k$). If $\eps=0$,
the shared-keep arm is therefore fresh-state Adam with the multiplier $\rho_k=1/\kappa(k,t{+}k)$. The
same holds for \texttt{reset\_v}, with the whole-update multiplier $\sqrt{b_{t+k}/b_k}$. These
identities describe the algorithms. Whether the empirical gain of an artifact arm is explained by its
schedule would require a learning-rate-matched run, which we did not perform.
\emph{(iii) What remains.} Under $\eps=0$, counter policies differ from one another only by learning-rate
schedules. Anything beyond scale must come from changing moment \emph{values}, as in \texttt{reset\_m},
\texttt{reset\_v} and \texttt{reset\_both}. This is what the update-norm controls test.

\paragraph{Exact correction does not bound the update.}
After a decoupled \texttt{reset\_v}, the first step has $\hat v_i=g_i^2$ exactly, while the kept first
moment gives $\hat m_i=(\beta_1 m_i^{\text{old}}+(1-\beta_1)g_i)/a_{t+1}$. Hence
\begin{equation}
|u_i| \;=\; \alpha\,\frac{|\hat m_i|}{|g_i|+\eps},
\label{eq:resetv}
\end{equation}
which is unbounded as $|g_i|/|\hat m_i|\to0$ and is capped only by $\alpha|\hat m_i|/\eps$. The
correction is exact, since $\hat v$ is unbiased for a stationary $g^2$, but $\hat v$ now rests on a single
sample. The familiar bound $|u_i|\lesssim\alpha(1-\beta_1)/\sqrt{1-\beta_2}$ \citep{kingma2015adam}
presumes that both moments start from zero together, and it no longer applies. Accuracy of the bias
correction and stability of the update are separate properties. For example, coordinates of a unit whose
current gradient is (near) zero but whose old momentum is not would receive very large steps. Our logs
store only per-tensor norms, so we cannot attribute the observed blow-up to particular coordinates
(\textsc{checkpoint\_trace\_unavailable}; Section~\ref{sec:interventions}).
